\documentclass[11pt,a4paper]{article}
\usepackage[utf8]{inputenc}
\usepackage[T1]{fontenc}
\usepackage[portuguese]{babel}
\usepackage{longtable}
\usepackage{booktabs}
\usepackage{array}
\usepackage{geometry}
\geometry{margin=2cm}
\usepackage{ragged2e}

\title{Dicionário de Dados --- IPDM da Baixada Santista}
\author{Projeto IPDM --- Ivan, Gustavo e Jorge}
\date{\today}

\begin{document}
\maketitle

\section*{Apresentação}
Este dicionário descreve as variáveis do conjunto de dados \texttt{df\_ipdm\_baixada\_por\_municipio.csv}, que reúne os índices sintéticos e componentes do Índice Paulista de Desenvolvimento Municipal (IPDM) para os nove municípios da Região Metropolitana da Baixada Santista, nos anos de 2014, 2016, 2018, 2020, 2022 e 2024.

\section*{Dicionário de Variáveis}

\begin{longtable}{@{}p{4.2cm} p{5.0cm} p{1.8cm} p{2.4cm} p{2.0cm}@{}}
\caption{Dicionário de dados do \texttt{df\_ipdm\_baixada\_por\_municipio}.}\\
\toprule
\textbf{Variável} & \textbf{Descrição} & \textbf{Tipo} & \textbf{Valores} & \textbf{Unidade de medida} \\
\midrule
\endfirsthead

\toprule
\textbf{Variável} & \textbf{Descrição} & \textbf{Tipo} & \textbf{Valores} & \textbf{Unidade de medida} \\
\midrule
\endhead

\midrule
\multicolumn{5}{r}{\small\textit{Continua na próxima página\ldots}}\\
\endfoot

\bottomrule
\endlastfoot

\texttt{cod\_ibge} &
Código oficial de 7 dígitos do município, definido pelo IBGE. &
Numérico (inteiro) &
7 dígitos (ex.: 3506359) &
--- \\

\texttt{Municipio} &
Nome do município pertencente à Região Metropolitana da Baixada Santista. &
Texto &
Bertioga, Cubatão, Guarujá, Itanhaém, Mongaguá, Peruíbe, Praia Grande, Santos, São Vicente &
--- \\

\texttt{Ano} &
Ano de referência do triênio em que o IPDM foi calculado. &
Numérico (inteiro) &
2014, 2016, 2018, 2020, 2022, 2024 &
Ano \\

\texttt{Riqueza [índice 0-1]} &
Indicador sintético da dimensão Riqueza do IPDM, composto por PIB per capita, rendimento do trabalho formal, e consumo de energia elétrica. &
Numérico (decimal) &
0,321 -- 0,575 &
Índice (0 a 1) \\

\texttt{Longevidade [índice 0-1]} &
Indicador sintético da dimensão Longevidade do IPDM, composto pelas taxas de mortalidade infantil, perinatal, de 15 a 39 anos e de 60 a 69 anos. &
Numérico (decimal) &
0,560 -- 0,756 &
Índice (0 a 1) \\

\texttt{Escolaridade [índice 0-1]} &
Indicador sintético da dimensão Escolaridade do IPDM, composto por atendimento escolar (0--3), proficiência no 5º e 9º ano e distorção idade-série. &
Numérico (decimal) &
0,334 -- 0,602 &
Índice (0 a 1) \\

\texttt{IPDM [índice 0-1]} &
Índice Paulista de Desenvolvimento Municipal: média aritmética das dimensões Riqueza, Longevidade e Escolaridade. &
Numérico (decimal) &
0,442 -- 0,604 &
Índice (0 a 1) \\

\texttt{Taxas de mortalidade 60 a 69 anos} &
Taxa de mortalidade da população de 60 a 69 anos, por mil habitantes. &
Numérico (decimal) &
13,94 -- 32,99 &
Por mil habitantes \\

\texttt{Taxas de distorção idade-série no Ensino Médio} &
Percentual de alunos do Ensino Médio com atraso escolar de dois anos ou mais. &
Numérico (decimal) &
8,20 -- 26,80 &
\% \\

\texttt{Produto Interno Bruto per capita} &
PIB per capita do município, em valores de 2024. &
Numérico (decimal) &
22.112 -- 270.118 &
R\$ (2024) \\

\texttt{Taxas de mortalidade 15 a 39 anos} &
Taxa de mortalidade da população de 15 a 39 anos, por mil habitantes. &
Numérico (decimal) &
1,07 -- 2,18 &
Por mil habitantes \\

\texttt{Rendimento do trabalho formal mais benefícios previdenciários per capita} &
Rendimento médio do trabalho formal somado a benefícios previdenciários (aposentadorias e pensões), per capita. &
Numérico (decimal) &
777,74 -- 3.183,93 &
R\$ (2024) \\

\texttt{Proporção média de alunos do 9º ano com proficiência} &
Média das proporções de alunos do 9º ano do EF da rede pública com proficiência adequada em Língua Portuguesa e Matemática (Prova Brasil). &
Numérico (decimal) &
12,74 -- 37,29 &
\% \\

\texttt{Taxas de mortalidade perinatal} &
Taxa de mortalidade perinatal (natimortos + óbitos de 0 a 6 dias), por mil nascidos vivos. &
Numérico (decimal) &
8,54 -- 21,38 &
Por mil nascidos vivos \\

\texttt{Proporção média de alunos do 5º ano com proficiência} &
Média das proporções de alunos do 5º ano do EF da rede pública com proficiência adequada em Língua Portuguesa e Matemática (Prova Brasil). &
Numérico (decimal) &
39,71 -- 71,87 &
\% \\

\texttt{Consumo anual de energia elétrica comercial, serviços e rural por ligação} &
Consumo médio anual de energia elétrica nas atividades comercial, de serviços e rural, por ligação elétrica. &
Numérico (decimal) &
3,08 -- 43,34 &
MWh por ligação \\

\texttt{Taxas de mortalidade infantil} &
Taxa de mortalidade infantil (óbitos de menores de 1 ano), por mil nascidos vivos. &
Numérico (decimal) &
7,80 -- 20,14 &
Por mil nascidos vivos \\

\texttt{Taxas de atendimento escolar de crianças de 0 a 3 anos} &
Percentual de crianças de 0 a 3 anos matriculadas em creche/escola em relação à população dessa faixa. &
Numérico (decimal) &
23,85 -- 58,11 &
\% \\

\texttt{Consumo anual de energia elétrica residencial por ligação} &
Consumo médio anual de energia elétrica residencial, por ligação elétrica. &
Numérico (decimal) &
2,43 -- 7,54 &
MWh por ligação \\

\end{longtable}

\section*{Observações}
\begin{itemize}
  \item Os valores em decimais usam vírgula como separador decimal, conforme o padrão brasileiro (ex.: \texttt{0,365}).
  \item Os indicadores sintéticos (\texttt{Riqueza}, \texttt{Longevidade}, \texttt{Escolaridade} e \texttt{IPDM}) variam entre 0 e 1.
  \item As faixas de valores apresentadas referem-se apenas aos nove municípios da Baixada Santista contemplados no recorte deste projeto, e não ao universo completo dos 645 municípios paulistas.
  \item Fonte original dos dados: Fundação Seade (IPDM) e demais órgãos citados no Anexo Metodológico (IBGE, MEC/Inep, Rais, INSS).
\end{itemize}

\end{document}